\begin{pastexam}[label={px:nsysu-111-dmalg-01}]{中山資工所資安組 111 學年度 第 2 題}
（計算）

{[10\%]} In how many ways can Alice place 36 different books on five shelves so that there is at least
two books on each shelf? (For any of these arrangements consider the books on each shelf to be placed
one next to the other, with the first book at the left of the shelf.)
\end{pastexam}

\begin{pastsoln}
答案：$\dbinom{30}{4} \cdot 36! = 27405 \cdot 36!$ 種。

書是不同的，書在架上的左右順序也算數，所以一種擺法由兩件事決定：五個書架各放幾本，以及書的排列。
把它拆成兩步：
\begin{enumerate}
\item \textbf{先把 36 本書排成一列}：$36!$ 種。
\item \textbf{再把這一列切成五段}，第 $i$ 段由左到右放到第 $i$ 個書架上。 設第 $i$ 個書架放 $x_i$
本，要求
\[
x_1 + x_2 + x_3 + x_4 + x_5 = 36, \qquad x_i \ge 2 .
\]
令 $y_i = x_i - 2 \ge 0$，變成 $y_1 + \cdots + y_5 = 36 - 10 = 26$ 的非負整數解，用隔板法
（$26$ 個球、$4$ 根隔板排成一列，資結書附錄~A.3~節）共有
\[
\binom{26 + 4}{4} = \binom{30}{4} = \frac{30 \cdot 29 \cdot 28 \cdot 27}{24} = 27405
\]
種。
\end{enumerate}
反過來，給一種擺法，把五個書架上的書依書架順序、由左到右接起來，就得到唯一的一個排列和唯一的
一組 $(x_1, \ldots, x_5)$。 所以「擺法」和「（排列，切法）」一一對應，總數是兩步相乘：
\[
27405 \cdot 36! .
\]
（用小例子檢查同樣的算法：$6$ 本書放 $2$ 個書架、每架至少 $2$ 本，公式給 $6! \cdot \binom{3}{1}
= 2160$；直接數，三種分法 $(2,4), (3,3), (4,2)$ 各有 $6!$ 種，合計 $2160$，一致。）
\end{pastsoln}
